% titel: Blatt 3, A4: Gradient eines Skalarfeldes
% skript: Kap. Vektoranalysis, Gradient und Rotation

% Gradient, Kraft = - nabla U, Rotation konservativer Felder.

Die Gravitationskraft ist der (negative) Gradient des Skalarpotentials
$U$; ein konservatives Feld hat verschwindende Rotation.

\begin{werkzeugkasten}[Werkzeuge]
\textbf{Gradient} (Vektor der partiellen Ableitungen):
\begin{equation}
  \nabla U=\left(\frac{\partial U}{\partial x},
  \frac{\partial U}{\partial y},\frac{\partial U}{\partial z}\right).
\end{equation}

\textbf{Kraft aus Potential:}
\begin{equation}
  \vec F=-\nabla U.
\end{equation}

\textbf{Radialsymmetrisches Feld:} die Kraft hat Betrag $|\vec F|$ und
zeigt in Richtung $-\hat e_r$ (Attraktion), falls $U$ mit wachsendem $r$
zunimmt (hier: $U\to 0$ fuer $r\to\infty$).

\textbf{Kettenregel fuer $r=\sqrt{x^2+y^2+z^2}$:}
\begin{equation}
  \frac{\partial U}{\partial x}=\frac{\mathrm dU}{\mathrm dr}
  \frac{\partial r}{\partial x},\qquad
  \frac{\partial r}{\partial x}=\frac{x}{r}\;.
\end{equation}

\textbf{Rotation:}
\begin{equation}
  \nabla\times\vec F=
  \begin{pmatrix}\partial_yF_z-\partial_zF_y\\
  \partial_zF_x-\partial_xF_z\\
  \partial_xF_y-\partial_yF_x\end{pmatrix}\;.
\end{equation}
Fuer ein Feld $\vec F=-\nabla U$ gilt stets $\nabla\times\vec F=\vec 0$;
eine verschwindende Rotation (wirbelfrei) ist die lokale Bedingung fuer
ein konservatives Kraftfeld.
\end{werkzeugkasten}

\begin{werkzeugkasten}[Vorgehen]
\begin{itemize}
  \item Zuerst $\partial r/\partial x_i=x_i/r$ bilden.
  \item $\nabla U$ als Vektor schreiben; das negative Vorzeichen durch
        $\vec F=-\nabla U$ macht die Kraft zum Zentrum hin.
  \item Fuer die Rotation eine Komponente exemplarisch pruefen oder die
        Identitaet ``Rotation eines Gradienten ist null'' ausnutzen.
  \item Mit dem Satz von Stokes ist die Wirbelfreiheit gleichbedeutend
        mit verschwindender Ringintegral.
\end{itemize}
\end{werkzeugkasten}

\begin{achtung}[Typische Fallen]
\begin{itemize}
  \item Beim Differenzieren von $U(r)$ die Kettenregel nicht vergessen
        ($r$ haengt von $x,y,z$ ab).
  \item Vorzeichen: $\vec F=-\nabla U$, nicht $+\nabla U$.
  \item $\nabla\times\vec F=\vec 0$ bedeutet: das Feld ist wirbelfrei und
        (in einfach zusammenhaengenden Gebieten) konservativ.
\end{itemize}
\end{achtung}
